\chapter{Adrenal Glands}
\index{Adrenal Glands}

\section*{Definition and Overview}
The adrenal glands are two small, triangular-shaped endocrine glands that sit on top of the kidneys, one on each side. They produce hormones that are essential for life and for the body's response to stress. Each gland has two distinct parts: the outer cortex, which produces cortisol, aldosterone, and small amounts of sex hormones, and the inner medulla, which produces the catecholamines adrenaline (epinephrine) and noradrenaline (norepinephrine).

Cortisol helps the body respond to stress, regulates metabolism and blood sugar, and dampens inflammation; aldosterone controls blood pressure and the balance of salt and water in the body; and the catecholamines drive the "fight or flight" response. Production of cortisol and aldosterone is controlled by the pituitary gland in the brain via a finely balanced feedback system. Disorders of the adrenal glands arise when they produce too much or too little hormone, or when a tumour grows within them, and all of these can have significant consequences for health.

\section*{Epidemiology}
Adrenal disorders are uncommon compared with many other endocrine conditions. Addison's disease (primary adrenal insufficiency) affects roughly 1 in 10,000 to 20,000 people and is somewhat more common in women, typically presenting in adulthood. Cushing's syndrome, caused by excess cortisol, is rare, but its most common cause -- long-term treatment with corticosteroid medicines -- is much more frequent in the population. Phaeochromocytomas and Conn's syndrome are both rare, although small, harmless adrenal lumps called incidentalomas are increasingly discovered by chance on abdominal scans.

Congenital adrenal hyperplasia, an inherited enzyme defect present from birth, affects roughly 1 in 10,000 to 18,000 births. The prevalence of both overproduction and underproduction of adrenal hormones tends to increase with age, and secondary adrenal insufficiency (arising from the pituitary suppressing the adrenal glands) is seen most often in people who have taken corticosteroids long term.

\section*{Aetiology and Pathophysiology}
Adrenal disease results from either too little or too much hormone production, or from growths within the glands.

\begin{itemize}
    \item \textbf{Adrenal insufficiency (Addison's disease):} destruction of the adrenal cortex, most often by autoimmune attack; other causes include tuberculosis, haemorrhage into the gland, and metastatic cancer
    \item \textbf{Secondary adrenal insufficiency:} suppression of the body's own cortisol production after prolonged use of corticosteroid medicines, or from pituitary disease
    \item \textbf{Cushing's syndrome:} excess cortisol, usually from long-term steroid medication, from a pituitary tumour driving the adrenals (Cushing's disease), or from a tumour within the adrenal gland itself
    \item \textbf{Phaeochromocytoma:} a usually benign tumour of the adrenal medulla that releases surges of adrenaline and noradrenaline
    \item \textbf{Conn's syndrome:} an adrenal adenoma producing excess aldosterone, causing hypertension and low potassium
    \item \textbf{Congenital adrenal hyperplasia:} an inherited enzyme deficiency that disrupts cortisol production and can cause excess androgen production
\end{itemize}

\section*{Clinical Features}
The clinical features depend on whether the glands are underactive or overactive.

\begin{itemize}
    \item \textbf{Underactive (Addison's disease):} profound fatigue, weight loss, low blood pressure with dizziness on standing, darkening of the skin, nausea, abdominal pain, salt craving, and low blood sugar
    \item \textbf{Excess cortisol (Cushing's syndrome):} weight gain around the face and trunk, a round "moon" face, purple striae, easy bruising, hypertension, impaired glucose tolerance or diabetes, muscle weakness, and mood changes
    \item \textbf{Excess catecholamines (phaeochromocytoma):} episodic severe headache, sweating, palpitations, and paroxysmal high blood pressure
    \item \textbf{Excess aldosterone (Conn's syndrome):} hypertension that may be resistant to treatment, muscle cramps, thirst, and weakness from low potassium
\end{itemize}

\section*{Differential Diagnosis}
The symptoms of adrenal disorders overlap with those of many more common conditions, and distinguishing between them is an important part of assessment.

\begin{itemize}
    \item Other causes of fatigue, including hypothyroidism, depression, anaemia, and chronic illness
    \item Essential hypertension and chronic kidney disease, which must be distinguished from Conn's syndrome and phaeochromocytoma
    \item Type 2 diabetes and the metabolic syndrome, which can mimic the metabolic features of Cushing's syndrome
    \item Panic attacks, hyperthyroidism, and cardiac arrhythmias, which can resemble catecholamine excess
    \item Secondary adrenal insufficiency from recent or ongoing steroid use
\end{itemize}

\section*{When to Seek Medical Care}
See a doctor for persistent fatigue accompanied by weight loss, darkening of the skin, or dizziness on standing, and for hypertension that is difficult to control or associated with headaches, sweating, and palpitations. Anyone taking long-term corticosteroid medicines should be reviewed regularly and should never stop them abruptly.

\textbf{Contact your local emergency services for:}
\begin{itemize}
    \item Collapse, severe weakness, vomiting, or confusion in a person known to have adrenal insufficiency -- features of an adrenal crisis
    \item Very low blood pressure or signs of shock
    \item Severe headache, chest pain, or palpitations in someone with a known adrenal tumour
    \item Sudden, severe abdominal or back pain with faintness, which may indicate bleeding into an adrenal gland
\end{itemize}

\section*{Diagnosis and Investigations}
Diagnosis is made with hormone blood tests and imaging, usually coordinated by an endocrinologist.

\begin{itemize}
    \item \textbf{Cortisol and ACTH:} an early-morning blood cortisol and a short ACTH (Synacthen) stimulation test confirm or exclude adrenal insufficiency
    \item \textbf{Cortisol excess:} late-night salivary cortisol, 24-hour urinary free cortisol, or a dexamethasone suppression test for Cushing's syndrome
    \item \textbf{Aldosterone and renin:} paired blood levels, with serum potassium, to investigate Conn's syndrome
    \item \textbf{Metanephrines:} breakdown products of catecholamines measured in plasma or urine to detect a phaeochromocytoma
    \item \textbf{Imaging:} CT or MRI of the adrenal glands to characterise tumours and incidentalomas
    \item \textbf{Newborn screening:} a heel-prick test for congenital adrenal hyperplasia is performed in some countries
\end{itemize}

\section*{Management}
Treatment aims to correct the hormone imbalance and deal with any underlying tumour.

\begin{itemize}
    \item \textbf{Addison's disease:} lifelong replacement of cortisol (hydrocortisone or prednisolone) and often aldosterone (fludrocortisone), with extra doses during illness or injury according to written "sick day rules"
    \item \textbf{Secondary insufficiency:} cautious tapering of corticosteroids under medical supervision to allow the glands to recover
    \item \textbf{Cushing's syndrome:} treatment of the cause, usually surgical removal of a pituitary or adrenal tumour, or careful tapering of steroid medication
    \item \textbf{Phaeochromocytoma:} blood-pressure control with alpha-blockers followed by surgical removal of the tumour
    \item \textbf{Conn's syndrome:} surgical removal of the adenoma, or a mineralocorticoid antagonist such as spironolactone where surgery is unsuitable
    \item \textbf{Congenital adrenal hyperplasia:} lifelong steroid replacement with dose adjustment through growth and pregnancy
    \item \textbf{Adrenal crisis:} urgent treatment with intravenous steroids, fluids, and correction of low blood sugar and potassium in hospital
\end{itemize}

\section*{Complications}
\begin{itemize}
    \item Adrenal crisis -- a life-threatening emergency of very low blood pressure, low blood sugar, and electrolyte disturbance in untreated or poorly managed insufficiency
    \item Complications of long-term cortisol excess, including diabetes, hypertension, osteoporosis, and increased susceptibility to infection
    \item Cardiovascular complications of hypertension in Conn's syndrome and phaeochromocytoma, including stroke and heart attack
    \item Over- or under-replacement of hormones, each with its own adverse effects
    \item Surgical risks, including bleeding and infection after adrenal tumour removal
\end{itemize}

\section*{Prognosis and Outcomes}
Addison's disease is very well controlled with lifelong replacement therapy, and most affected people live normal lives with careful attention to dose adjustment during intercurrent illness. Cushing's syndrome usually improves markedly once the cause is treated, although hypertension and diabetes may persist in some people. Phaeochromocytoma is cured by surgery in the majority of cases, and Conn's syndrome usually resolves or substantially improves. The outlook is best when the disorder is diagnosed early and managed by a specialist endocrinologist.

\section*{Prevention and Self-Care}
\begin{itemize}
    \item \textbf{Steroid use:} use corticosteroid medicines at the lowest effective dose, and never stop them abruptly
    \item \textbf{Medical alert:} if you have adrenal insufficiency, wear a medical alert bracelet and carry a steroid emergency kit and written sick-day instructions
    \item \textbf{Sick day rules:} know how to double or adjust your steroid dose during illness, injury, or surgery, and seek medical advice early when unwell
    \item \textbf{Education:} make sure family members and close contacts can recognise the signs of an adrenal crisis
    \item \textbf{Follow-up:} keep regular appointments with your endocrinologist and have annual reviews of blood pressure, electrolytes, and bone health
    \item \textbf{Health habits:} maintain a healthy weight, blood pressure, and bone strength through diet, exercise, and calcium and vitamin D adequacy
\end{itemize}

\section*{Sources and Further Reading}
The following sources provide reliable, up-to-date information on the adrenal glands and their disorders.

\begin{itemize}
    \item MSD Manual. \textit{Overview of the adrenal glands and adrenal disorders.} https://www.msdmanuals.com
    \item Mayo Clinic. \textit{Addison's disease, Cushing syndrome, and adrenal gland tumours.} https://www.mayoclinic.org
    \item NHS. \textit{Conditions: Addison's disease and Cushing's syndrome.} https://www.nhs.uk/conditions/
    \item MedlinePlus. \textit{Adrenal gland disorders.} https://medlineplus.gov
    \item National Institute of Diabetes and Digestive and Kidney Diseases. \textit{Adrenal insufficiency and related conditions.} https://www.niddk.nih.gov
    \item World Health Organization. \textit{Endocrine and metabolic health resources.} https://www.who.int
\end{itemize}
